% Electron-ion (Buneman) two-stream instability: electron beam through stationary ions
% Builds on Overleaf as-is (pdfLaTeX + BibTeX), or locally with
% `latexmk twostream_electron_ion.tex`. Companion: twostream_electron_electron.tex.
% Run: twostream/runs/buneman (input buneman.toml). Missing figures render as placeholders.
% Audience: an undergraduate taking a first course in plasma physics.
\documentclass[10pt,twocolumn,letterpaper]{article}

\usepackage[letterpaper,margin=0.75in,columnsep=0.25in]{geometry}
\usepackage[T1]{fontenc}
\usepackage[utf8]{inputenc}
\usepackage{lmodern}
\usepackage{microtype}
\usepackage{amsmath,amssymb}
\usepackage{bm}
\usepackage{graphicx}
\usepackage{xcolor}
\usepackage{booktabs}
\usepackage{siunitx}
\usepackage[font=small,labelfont=bf]{caption}
\usepackage[numbers,sort&compress]{natbib}
\usepackage[colorlinks=true,allcolors=blue!50!black]{hyperref}

% ---- Notation -------------------------------------------------------------
\newcommand{\wpe}{\omega_{pe}}        % electron plasma frequency
\newcommand{\wpi}{\omega_{pi}}        % ion plasma frequency
\newcommand{\vz}{v_0}                 % electron beam drift speed
\newcommand{\ppc}{N_\mathrm{ppc}}     % particles per cell
\newcommand{\dd}{\mathrm{d}}

% ---- Figure placeholder: \figorbox{file}{width}{label shown while missing} --
\newcommand{\figorbox}[3]{%
  \IfFileExists{figures/#1}%
    {\includegraphics[width=#2]{figures/#1}}%
    {\fbox{\parbox[c][0.55#2][c]{\dimexpr#2-2\fboxsep-2\fboxrule\relax}{\centering\small\textit{Placeholder:}\\ #3\\ \texttt{figures/\detokenize{#1}}}}}}

% ---- Margin notes for drafting (set \drafttrue -> \draftfalse to hide) -----
\newif\ifdraft \drafttrue
\newcommand{\todo}[1]{\ifdraft{\color{red}\textbf{[TODO:} #1\textbf{]}}\fi}

\title{\vspace{-1.5em}The Electron--Ion (Buneman) Two-Stream Instability:\\
       an Electron Beam through Stationary Ions}
\author{Harry Wu}
\date{\today}

\begin{document}
\maketitle

\begin{abstract}
\noindent
An electron beam streaming through ions that are initially at rest is unstable: small density
ripples grow exponentially into a large electric wave. This is the Buneman, or electron--ion
two-stream, instability. Cold-fluid theory predicts that the fastest-growing wave has
$k\vz\approx\wpe$ and growth rate $\gamma_{\max}\approx\tfrac{\sqrt3}{2}(m/2M)^{1/3}\wpe$, about
$0.054\,\wpe$ for hydrogen. Unlike the instability between two electron beams, this wave also
oscillates, at a real frequency of about $0.036\,\wpe$. We reproduce the instability with the
particle-in-cell (PIC) code \textsc{Entity} and compare the measured growth rate and frequency
with theory.
\todo{one sentence on the headline result once the run is in.}
\end{abstract}

% ===========================================================================
\section{Introduction}
\label{sec:intro}

Picture a beam of electrons streaming at speed $\vz$ through a plasma whose ions are at rest.
This is what happens when a current flows through a plasma. Suppose a small ripple appears in
the ion density. Its electric field slows some electrons and speeds up others, so the
electrons bunch. The electron bunches in turn
pull on the ions and deepen the original ripple. Under the right conditions this feedback runs
away and the ripple grows exponentially, an \emph{instability}~\citep{Buneman1959,Chen2016}.

The ions are essential. If they were infinitely heavy and could not move, the electrons alone
would only support ordinary plasma oscillations carried along by the beam, which do not grow.
Because ions are heavy, they respond slowly, and the growth is much slower than when two
electron beams stream through each other (companion report). Even so, the growth rate scales as
the cube root of the mass ratio, $(m/M)^{1/3}$, not its square root, so it remains fast enough to matter in real current-carrying
plasmas, where it heats the electrons and acts as an effective resistance~\citep{Buneman1959}.

As in the electron--electron case, the cold, uniform version can be solved exactly, so it is a
good test of a particle-in-cell (PIC) simulation~\citep{BirdsallLangdon1991}. Here we check that
a PIC simulation reproduces both the growth rate and the oscillation frequency predicted by
theory.

% ===========================================================================
\section{Linear theory}
\label{sec:theory}

\paragraph{Dispersion relation.}
Electrons of density $n_0$ stream at $+\vz$ through ions of density $n_0$ and mass $M$ that start
at rest. Both species are cold. The electron and ion plasma frequencies are
\begin{equation}
  \wpe^2 = \frac{n_0e^2}{\varepsilon_0 m}, \qquad \wpi^2 = \frac{n_0e^2}{\varepsilon_0 M}
  = \frac{m}{M}\,\wpe^2 .
\end{equation}
The derivation is the same as for two electron beams: linearize each species' continuity
equation and Newton's law for a ripple $\propto e^{i(kx-\omega t)}$, and insert the density
ripples into Gauss's law. Each species contributes one term $\omega_{pj}^2/(\omega-kv_{0j})^2$,
where $v_{0j}$ is its drift, so~\citep{Chen2016}
\begin{equation}
  1 = \frac{\wpi^2}{\omega^2} + \frac{\wpe^2}{(\omega - k\vz)^2}.
  \label{eq:disp}
\end{equation}
If the ions could not move ($\wpi\to0$) this gives $\omega = k\vz\pm\wpe$, which is real, so
nothing grows. With mobile ions, Eq.~\eqref{eq:disp} is a quartic in $\omega$. For some $k$ two
of its roots are complex, $\omega=\omega_r\pm i\gamma$, and the root with $\gamma>0$ grows as
$e^{\gamma t}$ while it oscillates at $\omega_r$.

\paragraph{Approximate solution.}
The strongest growth occurs when the beam electrons see the wave at their own natural
frequency, $k\vz\approx\wpe$. Because ions are heavy, the wave there is slow, $|\omega|\ll\wpe$.
Setting $k\vz=\wpe$ and expanding the electron term to first order in $\omega/\wpe$,
\begin{equation}
  \frac{\wpe^2}{(\omega-\wpe)^2} = \frac{1}{(1-\omega/\wpe)^2} \approx 1 + \frac{2\omega}{\wpe}.
\end{equation}
Equation~\eqref{eq:disp} then becomes $0 = \wpi^2/\omega^2 + 2\omega/\wpe$, or
\begin{equation}
  \omega^3 = -\frac{m}{2M}\,\wpe^3 .
\end{equation}
The cube roots of $-1$ include $e^{i\pi/3}=\tfrac12+i\tfrac{\sqrt3}{2}$, so the growing root is
\begin{equation}
  \omega \approx \Bigl(\frac{m}{2M}\Bigr)^{1/3}\wpe\Bigl(\frac12 + i\,\frac{\sqrt3}{2}\Bigr).
  \label{eq:approx}
\end{equation}
For hydrogen, $M/m=1836$, so $(m/2M)^{1/3}=0.065$ and $\gamma_{\max}\approx0.056\,\wpe$. The
cube root is why the growth is not as slow as the ion response alone would suggest: ions on
their own oscillate at $\wpi=0.023\,\wpe$, yet the wave grows more than twice as fast.

\paragraph{Exact solution.}
Solving the quartic numerically for each $k$ gives $\gamma_{\max}=0.0544\,\wpe$ at
$k\vz=1.006\,\wpe$, with $\omega_r=0.036\,\wpe$. The approximation in Eq.~\eqref{eq:approx} is
therefore good to about 3\% in $\gamma$. Waves with $k\vz>(1+(m/M)^{1/3})^{3/2}\,\wpe=1.125\,\wpe$
do not grow. The wave's phase speed $\omega_r/k\approx0.035\,\vz$ is nearly zero, so in the ion
frame the wave is almost standing still. This is what lets the slow ions keep up with it. These
exact $\gamma(k)$ and $\omega_r(k)$ curves are our reference. As in the electron--electron case,
growth stops once the wave is strong enough to trap the electrons.

% ===========================================================================
\section{The particle-in-cell method}
\label{sec:pic}

A PIC code follows ``macroparticles'', each standing for many real particles with the same
charge-to-mass ratio, and stores the fields on a grid of spacing $\Delta x$, with $\ppc$
particles per cell~\citep{BirdsallLangdon1991,HockneyEastwood1988}. Each time step $\Delta t$
(1) gathers the field at each particle from nearby grid points, (2) pushes the particles with
Newton's law, (3) deposits their current onto the grid, and (4) updates the fields. With no
magnetic field the push is the leapfrog scheme, $v^{n+1/2}=v^{n-1/2}+(q/m)E^n\Delta t$ and
$x^{n+1}=x^n+v^{n+1/2}\Delta t$~\citep{Boris1970}. The deposited current conserves charge
exactly~\citep{Esirkepov2001,Umeda2003}, and in 1D Amp\`ere's law reduces to
$\partial_t E_x=-J_x/\varepsilon_0$. Using few particles makes the density noisy,
$\delta n/n\sim\ppc^{-1/2}$. This noise supplies the initial ripple, so none has to be added by
hand. The companion report describes these steps in more detail.

% ===========================================================================
\section{Simulation setup}
\label{sec:setup}

We use \textsc{Entity}~\citep{Entity} (v1.3.3, \texttt{streaming} problem generator, input
\texttt{buneman.toml}) in 1D with periodic boundaries. Code units have $c=1$ and a reference
skin depth $d_0=10$. Electrons and ions each have density $n_0$, so $\wpe=0.1$.

\paragraph{Species.}
Species 1 is the electron beam, with four-velocity $u_x=\gamma_L v_x/c=0.1$
($\vz\approx0.0995\,c$). At this speed the Lorentz factor $\gamma_L\approx1$, so the beam is
non-relativistic. The small relativistic correction lowers $\wpe$ by 0.75\%. Species 2 is the
ions, with $M/m=1836$, initially at rest but free to move.

\paragraph{The return current.}
In a real plasma the current carried by the beam is balanced by the rest of the circuit. In a
1D periodic box there is nothing to balance it. Because \textsc{Entity} advances the field with
Amp\`ere's law, a uniform beam current would build up a uniform electric field and make the
whole beam slosh back and forth at $\wpe$ instead of streaming through the ions. To cancel it,
species 3 is a positive species that moves with the beam and carries the opposite current. Its
mass ($10^8\,m$) is so large that the field cannot change its velocity, so it acts as a fixed
external current. Species 4 is an equally heavy negative species that never moves and keeps
species 3 neutral. Together they add only $\sim10^{-8}\wpe^2$ to Eq.~\eqref{eq:disp}, which is
negligible. The electron--electron problem does not need this, because its two beams carry
equal and opposite currents.

\paragraph{Parameters.}
The fastest-growing wave has $k_{\max}\approx1.01$ ($\lambda\approx6.2$). In a box of length
$L=100$ the allowed waves have $k=2\pi m/L$, and mode $m=16$ sits on the peak. Modes up to
$m=17$ are unstable. The growth time is $1/\gamma_{\max}\approx184$, and the run lasts
$t=4000$, about 22 $e$-foldings. The beams are nearly cold: the thermal speed
$v_{th}=10^{-3}c$ is only 1\% of $\vz$, though thermal smearing ($k v_{th}/\gamma\approx0.2$)
may lower $\gamma$ slightly. The electron Debye length $\lambda_{De}=v_{th}/\wpe=0.01$ is about
one cell. Table~\ref{tab:params} lists all parameters.
\todo{shape order; update if the run changes any values.}

\begin{table}[t]
  \centering
  \caption{Simulation parameters from \texttt{buneman.toml} in code units ($c=1$), with
  dimensionless ratios on the right.}
  \label{tab:params}
  \small
  \begin{tabular}{@{}lll@{}}
    \toprule
    Parameter & Value & Normalized \\
    \midrule
    Beam four-velocity $u_x$            & $0.1$               & $\vz\approx0.0995\,c$ \\
    Density $n_e=n_i$                   & $n_0$               & $\wpe=0.1$ \\
    Ion mass $M$                        & $1836\,m$           & $\wpi=0.023\,\wpe$ \\
    Temperature $T/mc^2$               & $10^{-6}$           & $v_{th}/\vz\approx0.01$ \\
    Box length $L$                      & 100                 & $L\wpe/\vz\approx100$ \\
    Cells                               & 10240               & $\Delta x/\lambda_{De}\approx1$ \\
    Time step (CFL $=0.5$)              & $\Delta t\approx0.0049$ & $\wpe\Delta t\approx5\times10^{-4}$ \\
    Particles per cell $\ppc$           & 16                  & \\
    Current filter passes               & 4                   & \\
    Return-current mass                 & $10^8\,m$           & \\
    Run time                            & 4000                & $\gamma_{\max}t\approx22$ \\
    Shape order                         & \todo{}             & \\
    \bottomrule
  \end{tabular}
\end{table}

\paragraph{Measuring $\gamma$ and $\omega_r$.}
We Fourier transform $E_x(x,t)$ in $x$ to get the complex amplitude $\tilde E_k(t)$ of each box
wave. Its magnitude grows as $e^{\gamma t}$, so $\ln|\tilde E_k|$ versus $t$ is a straight line
of slope $\gamma(k)$. Its phase advances at the rate $\omega_r(k)$, which the
electron--electron case does not have, since there $\omega_r=0$. We fit both between the noise
floor and saturation (Fig.~\ref{fig:growth}).

\begin{figure}[t]
  \centering
  \figorbox{ei_growth_mode.pdf}{\columnwidth}{$\ln|\tilde E_k(t)|$ for mode $m=16$, with fit}
  \caption{Amplitude of the fastest-growing wave ($m=16$) versus time. The dashed line is the
  straight-line fit during exponential growth; its slope is $\gamma$. \todo{caption details.}}
  \label{fig:growth}
\end{figure}

% ===========================================================================
\section{Results}
\label{sec:results}

Figure~\ref{fig:disp} compares the measured growth rate and frequency of each mode with the
exact roots of Eq.~\eqref{eq:disp}.
\todo{describe agreement: measured $\gamma_{\max}/\wpe$ vs.\ 0.0544, $\omega_r/\wpe$ vs.\
0.036, location of the peak, behavior near the cutoff $k\vz=1.125\,\wpe$; any effect of the
small beam temperature.}

\begin{figure}[t]
  \centering
  \figorbox{ei_dispersion.pdf}{\columnwidth}{$\gamma/\wpe$ and $\omega_r/\wpe$ vs.\ $k\vz/\wpe$: theory and measured}
  \caption{Growth rate $\gamma$ and real frequency $\omega_r$ of each box mode (points) compared
  with the exact roots of Eq.~\eqref{eq:disp} (lines). \todo{caption details.}}
  \label{fig:disp}
\end{figure}

% ===========================================================================
\section{Discussion and conclusions}
\label{sec:discussion}

\todo{Summarize: (i) how well PIC recovers $\gamma(k)$ and $\omega_r(k)$; (ii) comparison with
the electron--electron case (slower growth, nonzero $\omega_r$, role of ion inertia); (iii)
limitations, e.g.\ small thermal spread, the return-current species, and only box modes are
sampled.}

{\small\setlength{\bibsep}{1pt}
\bibliographystyle{unsrtnat}
\bibliography{refs}}

\end{document}
